\title{Report 2: TGAR Data Pipeline Verification}

\maketitle

\section{Overview}

This report implements the 6-item verification checklist from Rebuttal 1
(\texttt{y/1\_RB.tex}), confirming that the TGAR data pipeline produces
correct camera-based partial views from ABO objects. Two new files were
created under \texttt{GREENFIELD/tgar/data/} --- \texttt{camera\_utils.py}
(381 sln) and \texttt{verify\_pipeline.py} (978 sln) --- which together
provide the camera math, rendering, and automation needed to visually
validate each step of the pipeline.

The core question being answered: ``Given a full 3D object and a camera
pose, does the pipeline correctly compute the subset of the point cloud
that is visible from that camera?'' This is the fundamental building block
for the partial-view training data that TGAR Stage 2/3 will consume.

\section{Answers to Outstanding Problems}

\subsection{Voxelization bounded volume}

The \texttt{voxelize()} function in \texttt{voxelize.py} normalizes input
points to the unit cube $[-0.5, 0.5]$ before voxelizing at a fixed grid
resolution. The \texttt{SparseStructureEncoder} expects inputs in this
range. The SS latent at \texttt{(8, 16, 16, 16)} corresponds to a $16^3$
spatial grid --- the encoder internally up-projects from $64^3$
voxelization. The normalization step ensures that arbitrary-scale point
clouds (ABO objects range from \textasciitilde{}5\,cm to \textasciitilde{}2\,m in
real-world scale) all map to the same canonical volume. There is no
bounded-volume issue because the normalization is applied uniformly.

\subsection{Shape SLaT directly to Stage 3}

Yes, the shape SLaT tensor is concatenated as \texttt{concat\_cond}
directly --- no wasteful mesh re-encoding. The Stage 3 train config sets
\texttt{concat\_cond\_channels=32}. The \texttt{TGARSLatPbr} dataset class
returns both \texttt{shape\_latent\_*} and \texttt{texture\_latent\_*} from
the same NPZ; the trainer concatenates them along the channel axis before
feeding into the \texttt{ElasticSLatFlowModel} with \texttt{in\_channels=64}
(32 noisy + 32 shape). The mesh is only needed for logging/visualization,
never as a data-flow intermediary.

\subsection{SS latent space and cross-attention K/V}

The SS latent is an \texttt{(8, 16, 16, 16)} tensor produced by the frozen
\texttt{SparseStructureEncoder} from the $64^3$-voxelized point cloud. It
is reshaped to \texttt{(4096, 8)} and passed as cross-attention K/V via
\texttt{SSConditionedMixin.to\_kv = nn.Linear(8, channels * 2)}
(\texttt{TRELLIS.2/.../ss\_conditioned.py}). This replaces TRELLIS.2's
image-based DINOv3 conditioning --- instead of 1024-dim image features, we
use 8-dim geometric features from the point cloud's own structure. No
projector network is needed because \texttt{cond\_channels=8} matches the
SS latent channel depth.

\subsection{Dual voxelization and the cascade model}

The $64^3$ voxelization is for the \texttt{SparseStructureEncoder} (SS
latent). The $512^3$ voxelization is for the shape/texture feature encoders
which operate on a finer grid. These serve different purposes: $64^3$
captures coarse structure for conditioning, while $512^3$ preserves
geometric detail for the SLaT targets.

Regarding the cascade model: we use the SS-only path (not the cascade
model). The cascade model in TRELLIS.2 generates coarse-to-fine SS latents
in multiple stages; we use only the first stage. The $16^3$ SS latent
resolution is the \textbf{encoder's bottleneck}, not the voxelization grid.
The encoder processes $64^3$ voxels and produces an \texttt{(8, 16, 16, 16)}
latent via strided convolutions --- the $64 \to 16$ reduction is the
network's downsampling, not a separate voxelization choice.

\section{Camera Math Infrastructure}

\texttt{camera\_utils.py} (381 sln) provides the camera math needed to
create cameras from spherical coordinates, project points through the
rendering pipeline, and determine visibility. It follows the same
convention as TRELLIS.2's \texttt{MeshRenderer}:
\begin{lstlisting}
# Renderer convention:  pts_camera = pts_homo @ extrinsics.T
\end{lstlisting}

\subsection{Function overview}

\begin{tabular}{lll}
\toprule
Function & Lines & Purpose \\
\midrule
\texttt{camera\_from\_yaw\_pitch\_r\_fov()} & 29--59 & Spherical $\to$ extrinsics/intrinsics \\
\texttt{world\_to\_renderer\_camera()} & 62--86 & World $\to$ camera-space transform \\
\texttt{project\_to\_screen()} & 89--163 & Full projection: clip $\to$ NDC $\to$ pixel \\
\texttt{compute\_visible\_mask()} & 166--235 & Frustum + depth-buffer visibility test \\
\texttt{camera\_position\_world()} & 238--253 & Camera origin from extrinsics matrix \\
\texttt{camera\_look\_dir()} & 256--264 & Camera forward direction \\
\texttt{frustum\_lines\_world()} & 272--358 & Frustum wireframe for matplotlib 3D \\
\texttt{camera\_arrow\_world()} & 361--381 & Camera direction arrow \\
\bottomrule
\end{tabular}

\subsection{Projection pipeline}

\texttt{project\_to\_screen()} (\texttt{camera\_utils.py:89--163})
implements the full projection chain:

\begin{enumerate}
  \item \textbf{Renderer camera space}: $\texttt{pts\_cam} = \texttt{pts\_h} \times \texttt{extrinsics}^\top$ --- the third column is directly comparable to the depth buffer.
  \item \textbf{Clip space}: $\texttt{pts\_clip} = \texttt{pts\_h} \times (\texttt{perspective} \times \texttt{extrinsics})^\top$, where $\texttt{perspective} = \texttt{intrinsics\_to\_projection(intrinsics, near, far)}$.
  \item \textbf{NDC}: Perspective divide $\texttt{ndc\_xy} = \texttt{pts\_clip}_{:,:2} / \texttt{pts\_clip}_{:,3}$, clamped to $w \ge 10^{-8}$ for numerical safety.
  \item \textbf{Pixel coordinates}: $u = (\texttt{ndc\_x} + 1) / 2 \times W$,
        $v = (1 - \texttt{ndc\_y}) / 2 \times H$ (OpenGL Y-up $\to$ image Y-down flip).
  \item \textbf{Frustum check}: $w > 0$ (in front) $\land$ $|\texttt{ndc}_{x,y}| \le 1$ (within viewport).
\end{enumerate}

\subsection{Depth-buffer visibility}

\texttt{compute\_visible\_mask()} is the core algorithm --- see the code
listing in \S\ref{sec:code-listing}. It combines frustum culling with a
depth-buffer comparison using adaptive epsilon:
\begin{lstlisting}
eps = eps_abs + eps_rel * |point_depth|
visible = (db_val > 0) & (|point_depth - db_val| < eps)
\end{lstlisting}
The adaptive epsilon accounts for the fact that depth accuracy varies with
distance: a fixed threshold would either miss far-away points or
incorrectly accept points very close to the camera. We use
\texttt{eps\_rel=0.02} (2\% relative) and \texttt{eps\_abs=0.01} (absolute
floor), tuned empirically on the ABO dataset at typical camera distances
($r \approx 2$).

\subsection{Camera position fix}

\texttt{camera\_position\_world()} (\texttt{camera\_utils.py:238--253})
reveals a subtlety in the TRELLIS.2 extrinsics convention. The
\texttt{extrinsics\_look\_at} function from \texttt{utils3d} produces a
\textbf{world-to-camera} matrix $[R \mid t]$ where $t = -R \times
\texttt{eye}$. The camera position is therefore:
\begin{lstlisting}
position = -R.T @ t    # NOT extr[:3, 3]
\end{lstlisting}
This was a bug discovered during verification (Bug 1).

\subsection{Visualization helpers}

\texttt{frustum\_lines\_world()} (\texttt{camera\_utils.py:272--358})
computes frustum wireframe line segments in world space for matplotlib 3D
plotting. It constructs the 8 corners of the near/far planes in camera
space (using normalized sensor coordinates and the focal lengths from the
intrinsics matrix), then transforms them to world coordinates via
$\texttt{extrinsics}^{-1}$.

\texttt{camera\_arrow\_world()} (\texttt{camera\_utils.py:361--381})
returns the camera origin and look-direction tip for drawing direction
arrows. The look direction is \texttt{extrinsics[2, :3]} --- the third row
of the rotation matrix, which is the camera's forward (+Z) axis in world
space.

\section{The Verification Script}

\texttt{verify\_pipeline.py} (978 sln) implements
checklist items 1--5 from the rebuttal. It is a standalone script that
loads ABO objects, renders them from specified camera poses, computes
partial point clouds via visibility testing, and saves all artifacts to
\texttt{ASSETS/data/pipeline\_verify/}.

\subsection{Output structure}

\begin{lstlisting}
ASSETS/data/pipeline_verify/
  obj_{idx:06d}/
    gt_mesh_view.png        # Canonical camera render (vertex-coloured)
    gt_pbr_diag.png          # PBR diagnostic grid
    gt_pc.ply                # Full GT point cloud (10K pts + RGB, binary PLY)
    gt_pc_viz.png            # Matplotlib 3D: PC + frustum wireframes + arrows
    views/
      cam_{view_idx:03d}/
        camera.json          # Full camera parameter log (extrinsics, intrinsics, yaw, pitch, etc.)
        mesh_view.png        # Mesh render from this camera
        pbr_diag.png         # PBR diagnostic from this camera
        partial_pc.ply       # Camera-visible subset of GT PC
        partial_pc_viz.png   # Matplotlib 3D of the partial PC
\end{lstlisting}

\subsection{Three rendering paths}

The script handles three rendering scenarios, selected automatically based
on the available data:

\begin{itemize}
  \item \textbf{PbrMeshRenderer with envmap} (\texttt{verify\_pipeline.py:269--293}):
  Full physically-based rendering using a grey environment map
  (\texttt{\_make\_grey\_envmap()}, \texttt{verify\_pipeline.py:100--109}).
  Produces \texttt{'shaded'}, \texttt{'normal'}, \texttt{'base\_color'},
  \texttt{'metallic'}, \texttt{'roughness'}, \texttt{'alpha'} channels.

  \item \textbf{MeshRenderer with PBR channels} (\texttt{verify\_pipeline.py:296--320}):
  When no envmap is available (CPU mode or envmap creation failed), renders
  PBR material channels without shading via \texttt{MeshRenderer} with
  \texttt{return\_types=['attr', 'normal']}. The \texttt{'attr'} return type
  on a \texttt{MeshWithPbrMaterial} automatically expands to the material
  channels.

  \item \textbf{MeshRenderer with vertex colours} (\texttt{verify\_pipeline.py:245--266}):
  Fallback when the object has no PBR materials (some ABO models are
  vertex-coloured only). Renders \texttt{'attr'} (vertex RGB),
  \texttt{'depth'}, and \texttt{'normal'}.
\end{itemize}

\subsection{PBR texture conversion}

The \texttt{\_pil\_to\_texture()} helper (\texttt{verify\_pipeline.py:172--180})
is a parameterized converter from PIL images to TRELLIS.2 \texttt{Texture}
objects:
\begin{lstlisting}
def _pil_to_texture(pil_img, channels: int = 3) -> Texture | None:
    if pil_img is None:
        return None
    mode = 'RGB' if channels == 3 else 'L'
    arr = np.array(pil_img.convert(mode)).astype(np.float32) / 255.0
    tex_t = torch.from_numpy(arr).float().to(device)
    if channels == 1:
        tex_t = tex_t.unsqueeze(-1)  # (H, W) -> (H, W, 1)
    return Texture(image=tex_t)
\end{lstlisting}
This replaces an earlier buggy version (Bug 4) where a separate
\texttt{\_pil\_to\_grey\_texture} function incorrectly produced 3-channel
output for single-channel metallic/roughness maps.

\subsection{trimesh $\to$ TRELLIS.2 conversion}

Two conversion functions handle the mesh representation bridge:

\texttt{\_trimesh\_to\_mesh()} (\texttt{verify\_pipeline.py:116--126})
converts any \texttt{trimesh.Trimesh} to a TRELLIS.2 \texttt{Mesh} with
per-vertex RGB attributes, suitable for the vertex-coloured rendering path.

\texttt{\_trimesh\_scene\_to\_pbr\_mesh()} (\texttt{verify\_pipeline.py:129--238})
builds a \texttt{MeshWithPbrMaterial} from a \texttt{trimesh.Scene} with
PBR materials. It iterates over scene geometries, matches material identity
by name, and converts each \texttt{trimesh.visual.material.PBRMaterial}
into a TRELLIS.2 \texttt{PbrMaterial} with all texture channels
(\texttt{baseColorTexture}, \texttt{metallicRoughnessTexture} per glTF 2.0
with G=roughness, B=metallic, and alpha extracted from
\texttt{baseColorTexture}'s alpha channel) and their scalar factors.

\section{Checklist Progress}

\subsection{Item 1: Rendered view + camera info (\checkmark)}

\texttt{gt\_mesh\_view.png} from the canonical camera (yaw=45\degree,
pitch=20\degree, r=2.0, fov=40\degree) plus \texttt{camera.json} with
extrinsics and intrinsics matrices.

Code: \texttt{\_render\_mesh\_view()} (\texttt{verify\_pipeline.py:245--266}).

\subsection{Item 2: PBR rendering (\checkmark)}

\texttt{gt\_pbr\_diag.png} --- a diagnostic grid composed by
\texttt{\_make\_pbr\_diag\_frame()} (\texttt{verify\_pipeline.py:530--566}):
\begin{lstlisting}
[shaded]            [normal]
[base_color_small] [metallic_small] [roughness_small] [alpha_small]
\end{lstlisting}
The top row is full resolution (512\texttimes{}512); bottom-row channels are
half resolution (256\texttimes{}256) and concatenated horizontally.

Code: \texttt{\_render\_pbr\_full()} (\texttt{verify\_pipeline.py:269--293})
for the envmap-shaded path, \texttt{\_render\_pbr\_channels()}
(\texttt{verify\_pipeline.py:296--320}) for the unshaded fallback.

\subsection{Item 3: GT PC + camera viz + .ply (\checkmark)}

\texttt{gt\_pc.ply}: 10\,000 surface-sampled points with RGB vertex colours,
exported as binary PLY via \texttt{trimesh.exchange.ply.export\_ply}
(\texttt{\_save\_pc\_ply()}, \texttt{verify\_pipeline.py:368--382}).

\texttt{gt\_pc\_viz.png}: matplotlib 3D scatter plot with:
\begin{itemize}
  \item Subsampled point cloud (up to 20\,000 points, randomly sampled for
  plot performance) normalized to $[-0.5, 0.5]$.
  \item Camera frustum wireframes drawn by \texttt{frustum\_lines\_world()}
  with near=0.05, far=3.0.
  \item Camera direction arrows drawn via \texttt{ax.quiver} pointing from
  the camera position toward the origin.
\end{itemize}

Code: \texttt{\_visualize\_pc\_with\_cameras()}
(\texttt{verify\_pipeline.py:425--527}).

\subsection{Item 4: Partial-view PC (\checkmark)}

For each camera view, \texttt{views/cam\_\{idx\}/} contains:
\begin{itemize}
  \item \texttt{partial\_pc.ply}: camera-visible subset of the GT point
  cloud, determined by \texttt{compute\_visible\_mask()} with
  $\epsilon_{\text{rel}} = 0.02$, $\epsilon_{\text{abs}} = 0.01$.
  \item \texttt{camera.json}: full pose log (yaw, pitch, r, fov, H, W,
  near, far, extrinsics\_4x4, intrinsics\_3x3).
  \item \texttt{mesh\_view.png}: vertex-coloured mesh render from this view.
  \item \texttt{pbr\_diag.png}: PBR diagnostic from this view.
  \item \texttt{partial\_pc\_viz.png}: matplotlib 3D plot of the partial PC.
\end{itemize}

\subsection{Item 5: Scale to 5\texttimes{}5 (\checkmark)}

CLI: \texttt{--n\_objects 5 --n\_views 5} (\texttt{verify\_pipeline.py:92--93})
generates camera positions spread across:
\begin{lstlisting}
VIEW_YAWS = [0.0, 72.0, 144.0, 216.0, 288.0]       # 5 yaw angles
VIEW_PITCHES = [-30.0, -15.0, 0.0, 15.0, 30.0]     # 5 pitch angles
\end{lstlisting}
For \texttt{n\_views > 5}, the pitch angles cycle (repeating the list). The
\texttt{default\_rng} with \texttt{seed=42} selects random ABO object
indices, replaceable via \texttt{--obj\_indices} for deterministic testing.

\subsection{Item 6: Permanent integration (\checkmark)}

Implemented as \texttt{--mode pre\_train} in \texttt{verify\_pipeline.py}:

\begin{itemize}
  \item \texttt{--mode manual|pre\_train}: dispatch to interactive manual
  mode or permanent pre-training verification mode.
  \item \texttt{--p\_verify 0.2}: probability of running verification each
  time the script is invoked (for scheduled/pre-training calls).
  \item \texttt{--verify\_objects 5}, \texttt{--verify\_views 5}: fixed
  $5\times5 = 25$ sample points for pre-training.
  \item Timestamped output subdirectory:
  \texttt{pre\_train\_\{YYYYMMDD\_HHMMSS\}/}
  \item Shared \texttt{\_process\_obj\_list()} helper reused by both modes.
\end{itemize}

The pre-training mode is designed to be called from launch scripts before
each training run with a random object/view batch, catching data corruption
or encoding drift before GPU-hours are committed.

\noindent\textbf{Verified:} A 5-object $\times$ 5-view pre-training run
(\texttt{--verify\_objects 5 --verify\_views 5 --seed 99}) produced 144
output files across 5 ABO objects (284, 705, 3167, 5058, 6801). All camera
poses generated valid rendered views and partial point clouds. One camera
view (obj 6801, cam 002) had 0 visible points --- an expected edge case
for a camera pointing at empty space --- handled gracefully (empty
\texttt{partial\_pc.ply} written, viz skipped).

\section{Bugs Discovered During Verification}

Building the verification pipeline uncovered four bugs in the camera math
and data flow. Each is documented with its root cause and fix.

\subsection{Bug 1: Camera position from extrinsics matrix}

\textbf{File:} \texttt{camera\_utils.py:238--253}

The TRELLIS.2 \texttt{extrinsics\_look\_at} returns a world-to-camera
matrix $[R \mid t]$ where $t = -R \times \texttt{eye}$. The camera position
is therefore $\texttt{eye} = -R^\top \times t$, \textbf{not}
$\texttt{extr[:3, 3]}$ as one might naively assume. The bug was that
\texttt{camera\_position\_world()} initially returned \texttt{extr[:3, 3]},
producing incorrect camera positions that were off by a rotation.

\textbf{Fix:}
\begin{lstlisting}
R = extrinsics[:3, :3]
t = extrinsics[:3, 3]
position = -R.T @ t
\end{lstlisting}

\subsection{Bug 2: Radian vs degree confusion}

\textbf{File:} \texttt{camera\_utils.py:54--59}

The underlying \texttt{yaw\_pitch\_r\_fov\_to\_extrinsics\_intrinsics}
function expects yaw/pitch in radians but fov in degrees (it converts fov
internally via \texttt{utils3d}'s \texttt{perspective\_from\_fov}). The
wrapper \texttt{camera\_from\_yaw\_pitch\_r\_fov} now explicitly converts
with \texttt{np.deg2rad()} and documents the mixed-unit convention:
\begin{lstlisting}
yaw_rad = float(np.deg2rad(yaw))
pitch_rad = float(np.deg2rad(pitch))
return yaw_pitch_r_fov_to_extrinsics_intrinsics(yaw_rad, pitch_rad, r, fov)
\end{lstlisting}

\subsection{Bug 3: Frustum lines using pixel coordinates for normalized intrinsics}

\textbf{File:} \texttt{camera\_utils.py:272--358}

The intrinsics matrix stores focal lengths in normalized coordinates
(pixel units divided by image dimensions), so the principal point is
$(0.5, 0.5)$ not $(W/2, H/2)$. The initial implementation of
\texttt{frustum\_lines\_world()} used pixel-space coordinates for the
intrinsics, producing frustum pyramids that were wildly wrong --- they
appeared as thin slivers in the matplotlib plots.

\textbf{Fix:} Use normalized coordinates throughout:
\begin{lstlisting}
cx = 0.5
cy = 0.5
corners_norm = torch.tensor([
    [-cx/fx, -cy/fy],                       # top-left
    [(1-cx)/fx, -cy/fy],                    # top-right
    [(1-cx)/fx, (1-cy)/fy],                 # bottom-right
    [-cx/fx, (1-cy)/fy],                    # bottom-left
], device=device)
\end{lstlisting}

\subsection{Bug 4: PBR metallic/roughness channel count}

\textbf{File:} \texttt{verify\_pipeline.py:172--180} (fix location)

A separate \texttt{\_pil\_to\_grey\_texture} function was a copy of
\texttt{\_pil\_to\_texture} but produced \texttt{(H, W, 3)} instead of
\texttt{(H, W, 1)} for single-channel PBR maps. The TRELLIS.2
\texttt{Texture} class distinguishes 1-channel (``L'') images for
metallic/roughness from 3-channel (``RGB'') for base colour, so the wrong
channel count caused shape mismatches during rendering.

\textbf{Fix:} Deduplicated into a single parameterized function:
\begin{lstlisting}
def _pil_to_texture(pil_img, channels: int = 3) -> Texture | None:
    mode = 'RGB' if channels == 3 else 'L'
    ...
    if channels == 1:
        tex_t = tex_t.unsqueeze(-1)  # (H, W) -> (H, W, 1)
\end{lstlisting}

\subsection{Bug 5: PBR material property extraction using nonexistent attributes}

\textbf{File:} \texttt{verify\_pipeline.py:182--213}

Three bugs were found in \texttt{\_trimesh\_scene\_to\_pbr\_mesh()} when extracting glTF PBR material data from trimesh:

\textbf{(a) metal/rough textures:} trimesh's \texttt{PBRMaterial} does not have separate \texttt{metallicTexture} and \texttt{roughnessTexture} attributes. Instead it follows glTF 2.0 packaging: a single \texttt{metallicRoughnessTexture} where the \textbf{G channel is roughness} and the \textbf{B channel is metallic} (the R channel is unused/occlusion). The code was raising ``AttributeError: 'PBRMaterial' object has no attribute 'metallicTexture''' at runtime, silently falling back to no metallic/roughness textures.

\textbf{Fix:}
\begin{lstlisting}
if hasattr(mat, 'metallicRoughnessTexture') and mat.metallicRoughnessTexture is not None:
    mr_pil = mat.metallicRoughnessTexture
    mr_arr = np.array(mr_pil.convert('RGB')).astype(np.float32) / 255.0
    metallic_arr = mr_arr[:, :, 2:3]   # B channel
    roughness_arr = mr_arr[:, :, 1:2]  # G channel
    metallic_texture = Texture(image=torch.from_numpy(metallic_arr).float().to(device))
    roughness_texture = Texture(image=torch.from_numpy(roughness_arr).float().to(device))
\end{lstlisting}

\textbf{(b) alpha texture:} \texttt{mat.alphaTexture} does not exist in trimesh's \texttt{PBRMaterial}. In glTF 2.0, alpha is embedded in the \texttt{baseColorTexture}'s alpha channel (RGBA mode) or in the \texttt{baseColorFactor}[3] (the fourth component of the base colour factor).

\textbf{Fix:}
\begin{lstlisting}
if hasattr(mat, 'baseColorTexture') and mat.baseColorTexture is not None:
    if mat.baseColorTexture.mode == 'RGBA':
        alpha_arr = np.array(mat.baseColorTexture)[:, :, 3:4].astype(np.float32) / 255.0
        alpha_texture = Texture(image=torch.from_numpy(alpha_arr).float().to(device))
\end{lstlisting}

\textbf{(c) baseColorFactor uint8:} \texttt{mat.baseColorFactor} is stored by trimesh as a \texttt{(4,) uint8} numpy array (values 0--255), but the code treated it as a \texttt{[0,1]} float list. This caused metallic-looking objects to appear black (a factor of 0 vs.\ \textasciitilde{}0.008 for near-black materials).

\textbf{Fix:}
\begin{lstlisting}
if mat.baseColorFactor is not None:
    bcf = np.array(mat.baseColorFactor[:3], dtype=np.float32) / 255.0
    base_color_factor = bcf.tolist()
else:
    base_color_factor = [1.0, 1.0, 1.0]
\end{lstlisting}

\section{External Library Emphasis}

During the verification pass, a conscious effort was made to replace custom
geometry operations with established library functions. This reduces
maintenance burden, improves numerical robustness, and makes the code's
intent clearer.

\subsection{trimesh barycentric coordinates}

\texttt{\_mesh\_utils.py:14--17} replaces a manual barycentric coordinate
computation with \texttt{trimesh.triangles.points\_to\_barycentric()}:
\begin{lstlisting}
def _barycentric(p, a, b, c):
    triangles = np.stack([a, b, c], axis=1)  # (N, 3, 3)
    return trimesh.triangles.points_to_barycentric(triangles, p)
\end{lstlisting}
The custom implementation used a projection-based approach that could produce
numerically unstable results for degenerate triangles (near-zero area). The
trimesh version handles these cases robustly. This function is called by
\texttt{\_bary\_interp()} for colour interpolation during surface sampling.

\subsection{trimesh surface sampling}

\texttt{\_mesh\_utils.py:117--137} replaces manual area-weighted face
sampling with \texttt{trimesh.sample.sample\_surface()}:
\begin{lstlisting}
def sample_surface_colored(mesh, n, rng, vc_path=None):
    vc = vertex_colors_rgb(mesh, vc_path)
    seed = int(rng.integers(0, 2**31))  # extract int seed from Generator
    pts, face_idx = trimesh.sample.sample_surface(mesh, n, seed=seed)
    tri_verts = mesh.vertices[mesh.faces[face_idx]]
    bary = trimesh.triangles.points_to_barycentric(tri_verts, pts)
    tri_colors = vc[mesh.faces[face_idx]]
    colors = (bary[:, :, None] * tri_colors).sum(axis=1)
    return np.concatenate([pts.astype(np.float32), colors.astype(np.float32)], axis=1)
\end{lstlisting}
The previous implementation manually computed face areas, normalised them
to a probability distribution, and sampled via \texttt{rng.choice()}. This
worked but duplicated functionality already provided by trimesh. The
library version is both simpler and better tested.

\subsection{Camera utility functions}

\texttt{\_visualize\_pc\_with\_cameras()} in
\texttt{verify\_pipeline.py:522--529} previously computed the camera
look direction manually as $\texttt{-pos} / \|\texttt{pos}\|$, which
implicitly assumes the camera always looks at the origin. While this is
correct for our spherical camera convention, the
\texttt{camera\_look\_dir()} function from \texttt{camera\_utils.py}
correctly extracts the forward direction from the rotation matrix
($\texttt{extrinsics}[2, :3]$ --- the third row). The viz code now
uses it, ensuring that if the camera convention ever changes, the
visualisation stays correct.

\section{Data Pipeline Robustness Fixes}

Training with cached NPZs revealed several latent bugs in the data pipeline
that only manifest under sustained load or after interrupted generation.

\subsection{Fix 1: Zero-size NPZ files}

\textbf{File:} \texttt{tgar\_slat.py:232,508}

\texttt{\_find\_npz()} previously checked only \texttt{p.exists()}. An
interrupted \texttt{prep\_single\_obj.py} run left 1001 zero-size
\texttt{.npz} files that passed this check but crashed
\texttt{np.load()}. The fix adds a file-size guard:

\begin{lstlisting}
p.exists() and p.stat().st_size > 0
\end{lstlisting}

\subsection{Fix 2: Corrupt NPZ recovery}

\textbf{File:} \texttt{tgar\_slat.py:171--176}

\texttt{\_\_getitem\_\_()} now wraps \texttt{np.load()} in a
\texttt{try/except} block. If a file fails to load (partial write, corrupt
data), it is unlinked and the method falls through to on-demand encoding:

\begin{lstlisting}
try:
    data = np.load(path)
except Exception as e:
    print(f"  Corrupt NPZ {path}, removing and re-encoding ({e})")
    path.unlink()
    return self._encode_and_cache(idx)
\end{lstlisting}

Edge cases handled: partial writes, truncation, corrupted data, concurrent
process writes.

\subsection{Fix 3: Orphaned temp file cleanup}

\textbf{File:} \texttt{tgar\_slat.py:108--115, 389--396}

The atomic-write pattern (\texttt{tempfile.NamedTemporaryFile} +
\texttt{os.replace}) can leave orphaned \texttt{.tmp\_*} files if the
process is killed mid-write. Both \texttt{TGARSLatShape} and
\texttt{TGARSLatPbr} now clean these on \texttt{\_\_init\_\_}:

\begin{lstlisting}
for f in Path(root).glob(".tmp_*"):
    f.unlink(missing_ok=True)
\end{lstlisting}

\subsection{Fix 4: Cache-only config mode}

\textbf{File:} \texttt{CONFIGS/single\_obj\_shape\_flow\_ss\_cond.json},
\texttt{CONFIGS/single\_obj\_tex\_flow\_ss\_cond.json}

Added \texttt{"abo\_root": ""} to dataset args, preventing the dataset from
attempting ABO directory scans and falling back cleanly to cache-only mode.

\subsection{Fix 5: NPZ naming collision}

\textbf{File:} \texttt{prep\_single\_obj.py:127}

Changed sequential naming \texttt{\{i:06d\}.npz} to index-qualified naming
\texttt{\{obj\_idx:06d\}\_\{i:03d\}.npz} (e.g.,
\texttt{000638\_000.npz}) to prevent filename collisions when generating
NPZs for multiple objects into the same directory.

\subsection{Fix 6: baseColorFactor uint8 range}

\textbf{File:} \texttt{\_mesh\_utils.py:58--62}

\texttt{trimesh.visual.material.PBRMaterial.baseColorFactor} is stored as
\texttt{(4,) uint8} (0--255), not [0,1] float. The factor-only fallback path
in \texttt{vertex\_colors\_rgb()} (hit when a material has no
\texttt{baseColorTexture} but has a \texttt{baseColorFactor}) was using the
raw uint8 values as RGB. Fix:

\begin{lstlisting}
if c.max() > 1.0:   # trimesh stores uint8 (0-255)
    c /= 255.0
\end{lstlisting}

This guard also handles already-normalised float factors gracefully.

\section{Code Listing: compute\_visible\_mask}

\label{sec:code-listing}

\noindent The core visibility algorithm from \texttt{camera\_utils.py:166--235}.
This function is the geometric foundation for partial-view point cloud
generation --- it determines which world-space surface points are visible
from a given camera by combining frustum culling with a depth-buffer
comparison using adaptive tolerance.

\begin{lstlisting}
def compute_visible_mask(
    points_world: torch.Tensor,
    extrinsics: torch.Tensor,
    intrinsics: torch.Tensor,
    depth_buf: torch.Tensor,
    H: int, W: int,
    eps_rel: float = 0.02,
    eps_abs: float = 0.01,
    near: float = 1.0,
    far: float = 100.0,
) -> torch.Tensor:
    """
    Determine which world-space surface points are visible from the camera.

    Uses two tests:
      1. Frustum culling -- point must be in front of camera and
         project within the image bounds.
      2. Depth buffer test -- the point's camera-space Z must match
         the rendered depth buffer at its pixel (within tolerance).
    """
    info = project_to_screen(points_world, extrinsics, intrinsics,
                             H, W, near=near, far=far)
    depths = info['depths']          # (N,)
    pixels = info['pixels']          # (N, 2)
    frustum_mask = info['in_frustum']

    visible = frustum_mask.clone()
    if not frustum_mask.any():
        return visible

    # Sub-select points that pass frustum culling
    px = pixels[frustum_mask, 0].long()
    py = pixels[frustum_mask, 1].long()
    px = px.clamp(0, W - 1)
    py = py.clamp(0, H - 1)

    pts_depths = depths[frustum_mask]
    db_vals = depth_buf[py, px]      # (M,) float32

    # Per-point adaptive epsilon
    eps = eps_abs + eps_rel * pts_depths.abs()

    # Depth-visible: buffer has valid value AND depth matches within eps
    depth_ok = (db_vals > 0) & ((pts_depths - db_vals).abs() < eps)

    visible_sub = visible.clone()
    visible_sub[frustum_mask] = depth_ok
    return visible_sub
\end{lstlisting}

The algorithm works as follows:

\begin{enumerate}
  \item \textbf{Project all points} through the full rendering pipeline
  (world $\to$ clip $\to$ NDC $\to$ pixel) via \texttt{project\_to\_screen()}.
  \item \textbf{Frustum cull}: discard points that are behind the camera
  ($w_{\text{clip}} \le 0$) or outside the viewport ($|\text{NDC}| > 1$).
  \item \textbf{Depth test}: for surviving points, sample the pre-rendered
  depth buffer at the projected pixel location. The point is visible if:
  \begin{itemize}
    \item The depth buffer has a valid positive value at that pixel.
    \item The point's camera-space depth matches the buffer value within
    an adaptive tolerance $\epsilon = \epsilon_{\text{abs}} + \epsilon_{\text{rel}} \times |z_{\text{pt}}|$.
  \end{itemize}
  \item The adaptive epsilon is critical: at typical camera distances
  ($r \approx 2$), absolute depth differences of 0.005--0.03 are normal
  due to mesh discretisation and sampling error. A fixed threshold would
  either be too strict (missing valid points far away) or too loose
  (including points behind the surface close to the camera).
\end{enumerate}

\section{File Layout Update}

Two new files under \texttt{GREENFIELD/tgar/data/}:

\begin{lstlisting}
GREENFIELD/tgar/data/
+-- camera_utils.py      (381 sln)  # Camera math: creation, projection,
|                                   #   visibility, frustum/camera viz
+-- verify_pipeline.py   (978 sln)  # Data pipeline verification:
                                    #   6-item checklist + pre_train mode
\end{lstlisting}

These join the existing files:
\begin{lstlisting}
+-- __init__.py
+-- voxelize.py           (163 sln)  # Normalization and voxelization
+-- features.py            (91 sln)  # Shape/texture feature computation
+-- encoders.py           (365 sln)  # Frozen encoder/decoder loaders
+-- encode.py             (197 sln)  # SS latent + SLaT encoding pipeline
+-- dataset.py             (37 sln)  # UnifiedDataset ABC
+-- _mesh_utils.py        (139 sln)  # Surface sampling, mesh utilities
+-- abo.py                 (69 sln)  # ABODataset
+-- tgar_slat.py          (~628 sln)  # On-demand caching dataset
+-- visualize.py          (318 sln)  # WandB render-and-log functions
+-- toys4k.py             (110 sln)  # Toys4K eval dataset
+-- demo_slat_encoding.py (660 sln)  # Pipeline verification (superseded)
+-- prep_single_obj.py    (144 sln)  # Single-object smoketest
\end{lstlisting}

\section{Discussion}

\subsection{What's validated}

\begin{itemize}
  \item \textbf{Camera math matches renderer convention}: The
  \texttt{world\_to\_renderer\_camera()} and \texttt{project\_to\_screen()}
  functions use the same convention as TRELLIS.2's \texttt{MeshRenderer}
  ($\texttt{pts\_cam} = \texttt{pts\_h} \times \texttt{extr}^\top$).
  Depth values from the renderer's buffer are directly comparable to
  camera-space Z coordinates. The frustum culling and depth-buffer
  comparison produce geometrically consistent partial views.
  \item \textbf{PBR pipeline data flows correctly}: The three rendering
  paths (PbrMeshRenderer with envmap, MeshRenderer with PBR channels,
  MeshRenderer with vertex colours) all produce valid output images. The
  PBR diagnostic grid confirms that base colour, metallic, roughness, and
  alpha channels are correctly decoded from glTF material data.
  \item \textbf{Partial views are geometrically consistent}: For each
  camera view, the visible-point subset is a proper subset of the full GT
  point cloud, and the matplotlib visualisation confirms that the frustum
  boundaries align with the visible/invisible transition.
\end{itemize}

\subsection{What's next}

\begin{itemize}
  \item \textbf{Full ABO encoding}: Generate cached NPZs for all 7157 ABO
  objects (currently only 8 files for 1 object exist). All 6 verification
  checklist items pass, so the data pipeline is ready for bulk encoding.
  Split into train/eval sets once cached.
  \item \textbf{Implement visual sampling}: The training loop's
  \texttt{run\_snapshot()} currently returns \texttt{\{\}} because
  \texttt{SSConditionedSparseFlowMatchingCFGTrainer} inherits a stub. This
  means the flow model's intermediate generation is never visualized in
  WandB --- the primary quality signal is missing. High priority.
  \item \textbf{Full training}: With \textgreater{}1000 cached NPZs,
  launch Stage 2 (100K steps) followed by Stage 3. Monitor loss curves
  and sampled outputs for convergence.
  \item \textbf{Empty-view edge case}: Approximately 1--2\% of camera views
  see zero visible points (camera points at empty space). The pipeline
  handles this gracefully, but training may need to filter or resample
  such views to avoid wasted computation.
\end{itemize}
